\lesson{8}{Sep 22 2021 Wen ()}{Honors Scientific Knowledge}{Unit 1}

Science is often considered the pursuit and acquisition of knowledge. Therefore, any actions that mimic this search for answers can look like science. However, the key difference between an explanation built with science and one built without it is the gaining of evidence through experimentation and the scientific method.

Sometimes a practice or belief claims to be science, but it does not follow the scientific method or is not proven reliable through experimentation. This is called pseudoscience. Let's learn why these examples of pseudoscience are not considered examples of science.

Scientists seek to understand the world around us and to learn how things work. To find answers, they investigate all sorts of natural phenomena and conduct experiments to collect data for analysis and interpretation. But to prove that their interpretations are valid, and in turn become scientific explanations of natural phenomena, their investigations must meet certain criteria.

\begin{itemize}
    \item Following Logic \\
    Using logic to interpret experimental data is important to the validity of a scientific explanation. For example, if data show that water evaporates with the addition of heat, it would not be proper logic to assume the water spontaneously evaporated.
    \item Peer Review \\
    Frequent examinations by scientists result in some ideas being refuted and replaced with other ideas. This frequent examination and testing makes scientific explanations more valid and durable over time.
    \item Global Access \\
    Scientific knowledge is constantly changing and developing because new observations or predictions are made every day. All scientific knowledge should be examined and re-examined using the steps of the scientific method to collect new empirical evidence.
    \item Rules of Evidence \\
    It is important that scientists share their investigations and conclusions with others so they can be tested and used by scientists all over the world.
\end{itemize}

\subsubsection*{Hidden Phenomena}

A scientific explanation should be observed, tested, or measured to be considered valid and reliable. Yet, there are some natural events that just can't be observed or tested using the technology available. For instance, in the 16th century Nicolaus Copernicus first proposed a mathematical heliocentric model. It theorized that the sun, not Earth, was the center of our solar system. The model was supported by empirical evidence collected from astronomers tracking objects in the night sky. But the model could not be completely confirmed until we sent probes and people into space centuries later.

Also, there are things in ancient history, before the written word, that scientists can only theorize about based on studies of items from that period. For example, the existence of dinosaurs is supported by numerous sources of their skeletal remains collected all around the world. Although no one has seen a living dinosaur, studies of these remains support that they once lived in vast numbers.

\subsubsection*{Theories and Hypothesis}

Within the body of scientific knowledge are two very important components: hypotheses and theories. Both types of scientific knowledge are developed using background research and can change over time using the scientific method. The difference is that scientific theories are well-tested hypotheses. They describe why things happen using empirical evidence collected and tested by multiple scientists using multiple kinds of investigations over a significant period. Hypotheses, on the other hand, are predictions or tentative explanations of phenomena that have not yet been fully tested or have been dis-proven.

\newpage
